% Local fix (Jordan factors and zero eigenvalues).
% Source: Papers/2011.12127/TN-Review-main.tex:433--441.
As documented in \cite{gap:cpgsv21_correlator_diagonalizable_expansion},
the pure-exponential formula in \cite[Section~II.B.3]{Cirac2021Matrix}
omits both polynomial Jordan factors and finite nilpotent transients, which
the following all-separation expansion restores.

\begin{theorem}[Exact binomial expansion of physical correlations]
    \label{thm:physical_correlator_binomial}
    \lean{MPSTensor.physicalConnectedCorrelator_exists_binomial_expansion}
    \leanok
    \uses{def:physical_connected_correlator}
    Let $T=\E_A-P$ be the complementary transfer map used in the centered
    physical correlator. There is a finite set $S$ of eigenvalues of $T$
    and coefficients $c_{\mu j}$ such that, at every separation $n\ge0$,
    \begin{align}
        C_{\mathrm{phys}}(X,Y;n)
         & =\sum_{\mu\in S}\sum_{j=0}^{k_\mu-1}
        \binom nj\mu^{\max(n-j,0)}c_{\mu j},
        \label{eq:physical_correlator_binomial}
    \end{align}
    where $k_\mu$ is the least exponent at which the generalized
    $\mu$-eigenspace stabilizes. Terms with $j>n$ have zero binomial
    coefficient; the displayed exponent uses truncated natural subtraction.
    This formula includes $n=0$ and $\mu=0$, with $0^0=1$.
\end{theorem}
\begin{proof}
    \leanok
    The generalized-eigenspace decomposition writes the centered inner
    insertion as a finite sum of components $v_\mu$ satisfying
    $(T-\mu I)^{k_\mu}v_\mu=0$. Expanding $T^n=((T-\mu I)+\mu I)^n$ on each
    component by the commuting binomial formula and applying the outer
    physical insertion and trace leaves only powers below $k_\mu$, while
    terms with $j>n$ carry a zero binomial coefficient. The coefficients
    $c_{\mu j}$ are the resulting fixed trace pairings with $(T-\mu I)^jv_\mu$.
\end{proof}

\begin{theorem}[Nonzero exponentials and the nilpotent transient]
    \label{thm:physical_correlator_jordan_transient}
    \lean{MPSTensor.physicalConnectedCorrelator_exists_jordan_expansion}
    \leanok
    \uses{def:physical_connected_correlator}
    The same correlator has the all-separation form
    \begin{align}
        C_{\mathrm{phys}}(X,Y;n)
         & =t(n)+\sum_{\mu\in S\setminus\{0\}}p_\mu(n)\mu^n,
        \label{eq:physical_correlator_jordan_transient}
    \end{align}
    with $\deg p_\mu<k_\mu$ and $t(n)=0$ for $n\ge k_0$.
\end{theorem}
\begin{proof}
    \leanok
    \uses{thm:physical_correlator_binomial}
    For $\mu\ne0$, the falling-factorial polynomial
    $p_\mu(x)=\sum_{j<k_\mu}c_{\mu j}\mu^{-j}
        x(x-1)\cdots(x-j+1)/j!$ has degree strictly below $k_\mu$, and
    $p_\mu(n)\mu^n$ reproduces the $\mu$-term of the binomial sum at every
    $n\ge0$, including $j>n$.
    The $\mu=0$ sum forms the separate transient $t(n)$: for $n\ge k_0$, every
    index $j<k_0$ satisfies $j<n$, so its factor $0^{n-j}$ vanishes and
    $t(n)=0$ (whereas writing $t(n)$ as a polynomial times $0^n$ would fail at
    $0<n<k_0$).
\end{proof}

\begin{lemma}[A common physical normal gauge]
    \label{lem:physical_correlation_normal_gauge}
    \lean{MPSTensor.IsNormalTensor.exists_physicalCorrelation_normalGauge}
    \leanok
    \uses{def:nt_cpsv, def:gauge_equiv, def:left_canonical_tensor,
        def:physical_left_right_correlator}
    A normal tensor $A$ with leading transfer eigenvalue one admits a
    trace-preserving gauge representative $B$
    with positive trace-one fixed state $\rho_B$, and positive normalized
    left/right fixed matrices $\rho^L,\rho^R$ for $A$. At every separation
    and for every pair of finite-support physical operators, the left-right
    connected contraction for $A$ equals that for $B,\Id,\rho_B$.
\end{lemma}
\begin{proof}
    \leanok
    \uses{lem:ldp_normal_gauge_reduction,thm:physical_correlator_gauge}
    For a gauge transformation $B^i=UA^iU^{-1}$, the matrices
    $\rho^L=U^\dagger U$ and $\rho^R=U^{-1}\rho_B U^{-\dagger}$ are positive
    and normalized by congruence and trace cyclicity, and gauge covariance
    equates the left-right connected contraction of $A$ with that of
    $B,\Id,\rho_B$. This single gauge choice underpins both the exponential
    decay theorem and the Jordan spectral expansion.
\end{proof}

\begin{theorem}[Jordan-aware physical spectral expansion for normal tensors]
    \label{thm:correlator_sum_exponentials}
    \lean{MPSTensor.physicalConnectedCorrelator_exists_normal_jordan_expansion}
    \lean{MPSTensor.IsNormalTensor.exists_physicalLeftRightCorrelation_jordan_expansion}
    \leanok
    \uses{def:nt_cpsv, def:physical_connected_correlator,
        def:physical_left_right_correlator}
    Let $A$ be normal with leading transfer eigenvalue one. There are positive
    normalized left/right fixed matrices for which every finite-support
    physical connected correlator has the form
    \begin{align}
        C_{\rho^L,\rho^R}^{A}(X,Y;n)
         & =t(n)+\sum_{\mu\in S}p_\mu(n)\mu^n,
        \label{eq:correlations_spectral_expansion}
    \end{align}
    at every $n\ge0$. Here $0\notin S$, each $\mu\in S$ is a transfer
    eigenvalue of $A$ with $|\mu|<1$, and $t(n)=0$ for $n\ge D^2$.
    A single canonical representative $B$ gives all Jordan indices:
    $\deg p_\mu$ is strictly below the stabilization index of the generalized
    $\mu$-eigenspace of its complementary transfer map. No diagonalizability,
    assumed expansion, or trace-preserving condition on the original $A$
    is required.
\end{theorem}
\begin{proof}
    \leanok
    \uses{thm:physical_correlator_jordan_transient,
        lem:physical_correlation_normal_gauge}
    In trace-preserving gauge, every nonzero complementary eigenvalue is a
    transfer eigenvalue with a traceless eigenvector. Normality and the
    peripheral-spectrum theorem for irreducible channels place all complementary
    eigenvalues strictly inside the unit disk, while generalized-eigenspace
    stabilization by the ambient dimension $D^2$ forces the nilpotent transient
    to vanish for $n\ge D^2$. Evaluating the preceding Jordan expansion in the
    common normal gauge and pulling the physical contraction back to $A$ yields
    the expansion for $A$; because gauge similarity preserves transfer
    eigenvalues, the frequencies $\mu\in S$ are transfer eigenvalues of $A$.
\end{proof}
